% =====================================================================
%  CHAPTER 17: 学校健康教育（对应大纲第十章）
%  参考PPT：第十七章 学校健康教育（罗江洪，赣南医科大学）
%  往届笔记：第十五章 学校健康教育与健康促进
% =====================================================================
\chapter{学校健康教育}
\setcounter{chapter}{17}
\chapterintro{
\textbf{【掌握/双横线】}学校健康教育的概念；健康素养的概念；学校健康教育的基本内容（5个领域、5个水平）；学校健康教育的主要理论与模式（理性行动理论和计划行为理论、阶段变化理论/TTM、社会网络和社会支持理论、组织改变理论、PRECEDE-PROCEED模式、健康行为生态学模型）。\newline
\textbf{【熟悉/单横线】}学校专题教育（预防艾滋病、毒品预防、预防烟草、学校性教育）；学校健康教育的教学方法（传统式、参与式、间接传播方法）；学校健康教育的组织实施（课堂教学、教师培训、学校支持性环境）；学校健康教育的评价类型（形成评价、过程评价、效果评价）和评价方法（观察法、访谈法、调查问卷）；RE-AIM评价框架。
}

% =====================================================================
%  PART A: 学校健康教育概述
% =====================================================================
\section{学校健康教育概述}

\subsection{学校健康教育与健康素养}

\begin{defbox}
\textbf{健康教育（health education）：}通过有计划、有组织、有系统的健康信息传播和行为干预，促使个人或群体掌握健康知识、树立健康观念，自觉采纳有益于健康的行为和生活方式的教育活动和过程。

\textbf{学校健康教育（school health education）：}以促进学生健康为核心的教育活动与过程，对象包括大、中、小学学生和幼儿园学龄前儿童。通过有计划、有组织、多种形式的教育教学活动，使学生掌握卫生保健知识，增强学生自我保健意识，养成科学、文明、健康的生活方式和行为习惯，达到预防疾病、增进身心健康、全面提高学生健康水平的目的。
\end{defbox}

\begin{keybox}
\textbf{健康素养（health literacy）：}个人获取、理解、处理基本的健康信息与服务，并运用这些信息和服务作出有利于维护和促进自身健康决策的能力。健康素养是健康的重要决定因素，提高全民健康素养水平是提高全民健康水平\imp{最根本、最经济、最有效}的措施之一。

\textbf{学生健康素养：}学生通过各种渠道获取健康信息和正确理解健康信息，并运用这些信息维护和促进自身健康的能力。学生健康素养普遍提高，防病意识和健康管理能力显著增强，体质健康水平明显提升是新时代学校健康教育的工作目标之一。
\end{keybox}

\subsection{学校健康教育基本内容}

\begin{keybox}
\textbf{中小学与普通高等学校健康教育的5个领域：}

\begin{table}[H]
\centering
\caption{学校健康教育的5个领域}
\begin{tabularx}{\textwidth}{lX}
\hline
\textbf{中小学} & \textbf{普通高等学校} \\
\hline
① 健康行为与生活方式 & ① 健康生活方式 \\
② 疾病预防 & ② 疾病预防 \\
③ 心理健康 & ③ 心理健康 \\
④ 生长发育与青春期保健 & ④ 性与生殖健康 \\
⑤ 安全应急与避险 & ⑤ 安全应急与避险 \\
\hline
\end{tabularx}
\end{table}

\textbf{5个水平（按学段划分）：}
\begin{enumerate}[leftmargin=*]
    \item 水平一（小学1--2年级）
    \item 水平二（小学3--4年级）
    \item 水平三（小学5--6年级）
    \item 水平四（初中7--9年级）
    \item 水平五（高中10--12年级）
\end{enumerate}
\end{keybox}

\subsection{学校健康教育的教学方法}

\begin{keybox}
\textbf{一、传统式健康教育方法（traditional approach）：}
指课堂讲授、讲座、示教等，能够帮助儿童青少年系统地掌握知识，引导他们树立正确的健康观念。包括：\imp{课堂讲授}、\imp{讲座}、\imp{示教}。

\textbf{二、参与式健康教育方法（participatory approach）：}
采用学生喜闻乐见的方式，激发学生主动参与的热情，让学生在主动参与和探索中学习健康知识，发展并形成健康意识和行为。学生为主体，教师更多地发挥组织协调作用。包括：
\begin{itemize}[leftmargin=*]
    \item \imp{小组讨论和案例分析}
    \item \imp{头脑风暴（brainstorming）}
    \item \imp{角色扮演、小品和游戏活动}
    \item \imp{辩论和演讲}
    \item \imp{同伴教育（peer education）}
\end{itemize}

\textbf{三、间接传播方法（indirect approach）：}
利用大众媒体和视听材料，为学生提供个性化学习的机会，进一步深化对健康知识的理解和应用，进而形成积极的健康行为和态度。包括：\imp{大众媒体}、\imp{视听手段}、\imp{网络系统性学习}。
\end{keybox}

\subsection{学校健康教育的组织与实施}

\begin{keybox}
\textbf{一、课堂教学：}
\begin{itemize}[leftmargin=*]
    \item \imp{教学形式}：专业课程设置、学科融合、综合实践。
    \item \imp{教学方法}：传统式、参与式、间接传播方法结合使用。
    \item \imp{教学要求}：健康教育倡导以\imp{生活技能}为基础；课程设置符合学生认知水平；教学方案明确目标及重难点；课程结束前及时小结升华内容。
\end{itemize}

\textbf{二、教师培训：}
健康教育师资培训可采取的形式：培训班、教研活动、教学交流、示范课、观摩课等。

\textbf{三、学校支持性环境：}
\begin{itemize}[leftmargin=*]
    \item \imp{物质环境}——桌椅设备、教室采光照明、课程安排、规章制度。
    \item \imp{人际环境}——师生之间、同学之间的尊重、支持和关怀关系。
\end{itemize}
\end{keybox}

% =====================================================================
%  PART B: 学校专题健康教育
% =====================================================================
\section{学校专题健康教育}

\subsection{预防艾滋病教育}

\begin{keybox}
\textbf{一、学校艾滋病（AIDS）预防教育的目标：}

\imp{总目标：}(1)了解预防AIDS相关知识；(2)培养健康的生活方式；(3)增强自我保护意识；(4)抵御HIV侵袭的能力。

\imp{分目标：}
\begin{itemize}[leftmargin=*]
    \item \textbf{初中阶段}——了解AIDS的基本知识、预防方法和措施，培养自我保护意识。
    \item \textbf{高中阶段}——进一步了解预防与控制AIDS相关知识，正确对待HIV感染者和患者，学会保护自己，培养对自己、他人及社会的责任感。
\end{itemize}

\textbf{二、学校AIDS预防教育的内容：}

\begin{table}[H]
\centering
\caption{学校艾滋病预防教育内容（按学段）}
\small
\begin{tabularx}{\textwidth}{>{\raggedright\arraybackslash}p{5.5cm}>{\raggedright\arraybackslash}p{5.5cm}}
\hline
\textbf{初中阶段（6课时）} & \textbf{高中阶段（4课时）} \\
\hline
AIDS基本知识 & 艾滋病的流行趋势 \\
AIDS对人类社会（重点在个人及家庭）的危害 & 艾滋病对社会、经济所带来的危害 \\
判断安全行为与不安全行为 & HIV感染者和AIDS患者的区别 \\
拒绝不安全行为的技巧 & AIDS的窗口期、潜伏期 \\
如何寻求帮助的途径和方法 & 吸毒与AIDS \\
预防AIDS教育相关的青春期生理和心理知识 & 无偿献血知识 \\
 & 预防AIDS的方法和措施 \\
 & 了解歧视对AIDS防治工作的影响 \\
 & 性道德与法制教育 \\
 & 中国预防控制AIDS的相关政策 \\
\hline
\end{tabularx}
\end{table}
\end{keybox}

\subsection{预防毒品教育}

\begin{keybox}
\textbf{一、毒品预防教育的目标：}

\imp{总目标：}在各学科渗透毒品预防教育的基础上，通过专题教育的形式，培养学生健康的生活情趣、毒品预防意识和社会责任感，掌握一些自我保护的方法，做"珍爱生命、拒绝毒品"的人。

\imp{分目标：}
\begin{itemize}[leftmargin=*]
    \item \textbf{小学}——了解毒品危害的简单知识，远离毒品危害。
    \item \textbf{初中}——了解有关禁毒的法律知识，拒绝毒品诱惑。
    \item \textbf{高中}——学会自我保护，培养禁毒意识和社会责任感，发现可疑情况能够及时报告。
\end{itemize}

\textbf{二、毒品预防专题教育内容（按学段）：}

\begin{itemize}[leftmargin=*]
    \item \imp{小学阶段（4课时）}：知道常见毒品的名称；初步了解毒品对个人和家庭的危害；知道一些不良生活习惯可能会导致吸毒；懂得一些自我保护的常识和简单方法。
    \item \imp{初中阶段（6课时）}：知道毒品的概念，能识别常见毒品名称；进一步了解毒品对个人和社会的危害；知道吸毒是违法行为，走私、贩卖、运输、制造毒品是犯罪行为，都要受到法律的惩处；学会一些拒绝毒品的方法，能够保护自己不受毒品侵害。
    \item \imp{高中阶段（4课时）}：以\imp{新型毒品的伪装性、迷惑性、危害性和防范知识}为毒品预防专题教育的重点；懂得选择毒品就是自我毁灭，学会向毒品说"不"；了解当前禁毒工作面临的形势，增强禁毒意识；培养社会责任感，参与学校、社区组织的禁毒宣传活动。
\end{itemize}
\end{keybox}

\subsection{青少年预防烟草教育}

\begin{keybox}
\textbf{一、烟草在青少年中的流行情况：}

中国青少年尝试吸烟率和现在吸烟率逐年上升，吸烟人群年轻化趋势明显。中国疾病预防控制中心《2023年中国青少年烟草调查报告》表明，13.7\%的中学生尝试吸烟，男生（19.1\%）高于女生（7.8\%）。初中生尝试吸烟率为11.2\%，高中生尝试吸烟率为17.1\%。66.8\%的尝试行为发生在13岁及之前，超过半数的尝试行为发生在小学阶段。

\textbf{二、青少年预防烟草健康教育的意义：}

与成年人相比，青少年大脑中与认知控制有关的大脑回路还没有发育完全，更容易对尼古丁上瘾并且更有可能成为终生吸烟者。青少年吸烟也会导致睡眠时长缩短，同时增加入睡潜伏期，进一步影响睡眠质量。

\textbf{三、电子烟（electronic cigarette）：}

指用于将电子烟烟液雾化为可吸入气溶胶供人抽吸等的电子传送系统。不单是传统烟草，"印象中不会成瘾的"电子烟中也含有大量尼古丁、大麻提取物，且电子烟中添加的保湿剂加热会产生羰基化合物等有害物质，对青少年的身体健康（如肝脏、心脑血管、免疫系统和神经发育等）和心理健康（如抑郁、自杀倾向等）产生一定的负面影响。近年来，电子烟凭借"戒烟神器"的概念风靡市场，通过社交媒体宣传和新颖且具有迷惑性的包装，吸引儿童少年成为其一大受众，甚至在许多国家儿童少年的使用率高于成人（WHO, 2024）。

观察性研究的荟萃分析显示，电子烟使用与戒烟成功率无直接关联。电子烟气溶胶中的丙烯醛、甲醛等含羰基成分具有潜在的致癌性。

\textbf{四、青少年预防烟草健康教育的内容：}
\begin{enumerate}[leftmargin=*]
    \item 科学引导青少年学生树立正确的健康观，牢固树立"自己是健康第一责任人"的观念，倡导青少年"拒绝第一支烟"。
    \item 充分利用爱国卫生月、世界无烟日等主题活动，以易于接受的方式宣传烟草烟雾的危害，促进形成青少年无烟环境。
    \item 将烟草危害和二手烟危害等控烟相关知识纳入中小学学生健康教育课程。
    \item 主动加强对电子烟危害的宣传教育，倡导青少年远离电子烟。
\end{enumerate}
\end{keybox}

\subsection{学校性教育}

\begin{keybox}
\textbf{学校性教育（school-based sexuality education）：}以学校为场所，以学生为主要教育对象进行的性科学教育活动。

\textbf{目标：}
\begin{enumerate}[leftmargin=*]
    \item 针对儿童青少年性发育需要，提供准确信息。
    \item 为儿童青少年探索性和社会关系相关的价值观、态度和规范提供机会。
    \item 促进儿童青少年性发育、性保护技能的获得。
    \item 鼓励儿童青少年承担责任并尊重他人的权利。
\end{enumerate}

\textbf{《国际性教育技术指导纲要》推荐的性教育8个核心概念：}
\begin{enumerate}[leftmargin=*]
    \item \imp{关系}——家庭；友谊、爱情和亲密关系。
    \item \imp{价值观、态度、文化与性}——宽容、包容及尊重；长期承诺及子女养育；价值观与性；人权与性；文化、社会与性。
    \item \imp{理解社会性别}——社会性别及其规范的社会建构；社会性别平等、刻板印象与偏见；基于社会性别的暴力。
    \item \imp{暴力与安全保障}——暴力；许可、隐私及身体完整性；信息与通信技术的安全使用。
    \item \imp{健康与福祉、技能}——社会规范和同伴对性行为的影响；决策；沟通、拒绝与协商技巧；媒介素养与性。
    \item \imp{人体与发育}——性与生殖解剖及生理；生殖；青春发育期；身体意向。
    \item \imp{性与性行为}——性与性的生命周期；性行为与性反应。
    \item \imp{性与生殖健康}——怀孕与避孕；艾滋病病毒和艾滋病的污名、关爱、治疗及支持；理解、认识与减少包括艾滋病病毒在内的性传播感染风险。
\end{enumerate}
\end{keybox}

\newpage

% =====================================================================
%  PART C: 学校健康教育主要理论与模式
% =====================================================================
\section{学校健康教育主要理论与模式}

\subsection{微观层面的理论与模式}

\subsubsection{（一）理性行动理论和计划行为理论}

\begin{keybox}
\textbf{理性行动理论（Theory of Reasoned Action, TRA）：}分析态度如何有意识地影响个体行为的理论。该理论认为，个体行为由\imp{行为意向}直接决定，而行为意向又由\imp{态度}和\imp{主观规范}共同决定。

\textbf{计划行为理论（Theory of Planned Behavior, TPB）：}在TRA基础上增加了\imp{感知行为控制（perceived behavioral control）}变量。TPB认为，行为意向受态度、主观规范和感知行为控制三者共同影响。TPB在预测方面优于TRA，但对于感知行为控制较低的学生TRA更优。

\textbf{应用：}TRA和TPB均为预测青少年吸烟行为的有效工具。研究表明，青少年对吸烟的态度、社会规范和感知行为控制能够显著预测吸烟意向，吸烟意向在一定程度上介导了态度、社会规范和感知行为控制对吸烟行为的独立影响。
\end{keybox}

\subsubsection{（二）阶段变化理论（跨理论模式/TTM）}

\begin{keybox}
\textbf{跨理论模式（Transtheoretical Model, TTM）/ 阶段变化模式：}关注行为变化的解释（而非仅仅关注行为变化）、明确行为变化的时间维度，指导人们通过阶段化的规划做出行为改变过程的模式。TTM的独特性在于其在行为改变中指定了时间维度，提出人们在改变行为的同时会经历各个阶段，整个过程可以从6个月到5年。

\textbf{行为变化的5个阶段：}

\begin{table}[H]
\centering
\caption{TTM行为变化阶段}
\begin{tabularx}{\textwidth}{l>{\raggedright\arraybackslash}p{4cm}X}
\hline
\textbf{阶段} & \textbf{时间定义} & \textbf{特点} \\
\hline
\imp{预设阶段（precontemplation）} & 未来6个月 & 在可预见的未来不考虑改变 \\
\imp{沉思阶段（contemplation）} & 未来1$\sim$6个月 & 考虑改变但未立即行动 \\
\imp{准备阶段（preparation）} & 下个月 & 计划在不久的将来进行改变 \\
\imp{行动阶段（action）} & 过去6个月中 & 已发生了有意义的变化 \\
\imp{维持阶段（maintenance）} & 6个月或更长时间 & 已将该变化保持了一段时间 \\
\hline
\end{tabularx}
\end{table}

\textbf{应用：}基于TTM对美国中学生酒精、烟草和其他药物使用情况的研究发现，网络个性化互动干预可显著增加"停止使用者"比例。
\end{keybox}

\subsection{宏观层面的理论与模式}

\subsubsection{（一）社会网络和社会支持理论}

\begin{keybox}
\textbf{社会网络（social network）：}社会成员之间因为互动而形成的相对稳定的关系体系，关注成员之间的互动和联系，其中成员可以是个体的，也可以是一个组织或家庭集合。

\textbf{核心发现：}青少年若其社交网络中吸烟者占比较高，则自身更易成为吸烟者。吸烟者社交圈中吸烟者比例较高，但网络大小、与网络成员接触频率及整体社会支持程度与非吸烟者相近。

\textbf{应用意义：}提示学校健康教育应重视\imp{同伴影响}和\imp{社交网络结构}在健康行为形成和改变中的作用，可以通过同伴教育、榜样示范等方式发挥社会网络的正向作用。
\end{keybox}

\subsubsection{（二）组织改变理论}

\begin{keybox}
\textbf{组织改变理论模型（Theory of Organizational Change）：}通过人群所在组织的改变、规章制度和管理模式的完善，实现对组织内全体人群的干预。

\textbf{组织改变的4个阶段：}
\begin{enumerate}[leftmargin=*]
    \item \imp{确立问题或认知阶段}——提高组织高层管理者对问题的认知（如提高学校领导对"无烟学校"的认知）。
    \item \imp{初期行动阶段}——组建协调小组（如组建控烟协调小组），制订项目规划和实施计划。
    \item \imp{执行阶段}——营造氛围、充实活动、编制教材等具体干预措施的实施。
    \item \imp{制度化阶段}——总结工作，将有效措施制度化，使新政策变为确定的、融入组织的新的目标价值。
\end{enumerate}

\textbf{应用——学校"戒烟干预"：}在学校环境中开展戒烟干预研究，遵循以上4个阶段，从提高管理层认知、组建小组、营造戒烟氛围和充实学生课余活动，到最终将有效措施制度化。
\end{keybox}

\subsection{项目层面的理论与模式}

\subsubsection{（一）PRECEDE-PROCEED模式}

\begin{keybox}
\textbf{PRECEDE-PROCEED模式：}由美国健康教育专家Lawrence W. Green和Marshall W. Kreuter共同提出，是目前世界上最常用的健康教育与健康促进模式。

\textbf{模式组成：}
\begin{itemize}[leftmargin=*]
    \item \imp{PRECEDE（第一部分）}——包括教育和生态评价中的\imp{易感因素（predisposing factors）}、\imp{促成因素（enabling factors）}和\imp{强化因素（reinforcing factors）}。在对目标人群进行社会学、流行病学、教育生态学以及政策等因素进行综合评价后，找出相关影响项目成败的促进与制约因素，规划与评价项目的启动、实施和完成。
    \item \imp{PROCEED（第二部分）}——教育和环境改变中的\imp{政策（policy）}、\imp{管理（regulatory）}和\imp{组织策略（organizational constructs）}。
\end{itemize}

\textbf{核心逻辑：}先进行需求评估（PRECEDE），再实施干预和政策支持（PROCEED）。

\textbf{应用：}利用PRECEDE-PROCEED模型对初中女生的饮食习惯进行干预，发现干预组在知识和态度等易感因素、促成因素和健康饮食强化因子得分等方面均有显著改善。
\end{keybox}

\subsubsection{（二）健康行为生态学模型}

\begin{keybox}
\textbf{健康行为生态学模型（Ecological Model of Health Behavior）：}1988年由Kenneth R. McLeroy等学者提出，将生态学理论应用于健康教育与健康促进，认为要同时关注个人和社会因素，除教育活动外，还包括倡导、组织改变、政策形成、环境改变等多种方法策略。

\textbf{核心观点：}环境是影响行为的主要因素，环境能够影响的不仅仅是某一个体，而是一类群体；健康行为发生、发展层次受多个水平的因素影响。

\textbf{核心原则（4条）：}
\begin{enumerate}[leftmargin=*]
    \item \imp{构建多级别、多层次的健康行为生态学干预模型}。
    \item \imp{形成多层次、交互协作的健康行为干预社会支持网络}。
    \item \imp{实现多级别、多层次社会网络间的交互作用}。
    \item \imp{生态学环境内健康行为干预对象的广泛性}。
\end{enumerate}

\textbf{核心内容：}
\begin{enumerate}[leftmargin=*]
    \item \imp{行为生成的多维度视角}——行为受个体、人际、组织、社区和公共政策等多个层面因素共同影响。
    \item \imp{系统相互关系}——各层面因素之间存在复杂的相互作用。
    \item \imp{综合干预策略}——需在多层面同时采取干预措施，才能有效改变行为。
    \item \imp{目标人群}——干预可针对个体、群体或整个生态系统。
\end{enumerate}

\textbf{应用：}针对小学生肥胖的整群随机临床试验，采用针对儿童个体（健康教育+身体活动+BMI监测）和针对儿童环境（学校政策支持+家庭参与）的双层干预策略。
\end{keybox}

\newpage

% =====================================================================
%  PART D: 学校健康教育评价
% =====================================================================
\section{学校健康教育评价}

\subsection{评价类型}

\begin{keybox}
\textbf{一、形成评价（formative evaluation）：}
在健康教育项目计划设计阶段进行的评价，主要包括：了解目前学生的健康知识态度和健康相关行为、健康状况、学校现有的资源、实施方案的科学性和可行性、传播材料和测量工具的预实验与完善等。

\imp{评价指标：}实施计划的科学性、政策的支持性、技术的可行性、目标人群的接受程度等。

\textbf{二、过程评价（process evaluation）：}
在健康教育项目实施过程中进行的评价，主要内容包括：
\begin{itemize}[leftmargin=*]
    \item 课堂教学有无课时、教案、考试。
    \item 有无专、兼职教师。
    \item 教师是否经过培训、教学方法如何。
    \item 健康教育活动有无活动方案、活动过程资料、教职员工和学生参与情况等。
\end{itemize}

\textbf{三、效果评价（effect evaluation）：}
在健康教育项目实施后对效果进行的评价，主要指标包括：学生\imp{知识}、\imp{态度}、\imp{行为}的变化等。
\end{keybox}

\subsection{评价调查方法}

\begin{keybox}
\textbf{一、观察法：}
可用于评估技能和行为目标，也可用于简单评估知识和态度目标。教师可以直接观察学生表现出的与健康教育相关的技能。例如：在预防艾滋病的健康教育中，让学生自己设计海报、明信片等以反映出学生的态度倾向，或举办辩论赛以评估学生有关艾滋病的态度。

\textbf{二、个人访谈、小组访谈和资料与案例分析：}
通过面对面沟通的形式，获取学生对知识、态度和技能的掌握情况。访谈者可以是老师、学生或培训过的访谈员，允许学生自由地表达自己的想法。例如：预防近视的健康教育中，要求学生阅读一个案例并识别其中的危险行为（知识），描述应怎样保护视力（技能），以及描述对案例中人物行为的感想（态度）。

\textbf{三、调查问卷：}
设置封闭式和开放式问题，可用于知识、态度、行为的评价。
\begin{itemize}[leftmargin=*]
    \item \imp{封闭式问题}——常见形式有判断题、单选题、多选题和匹配题等，可采用分级测量设计，如利克特量表（Likert scale）、博格达斯社会距离量表等。
    \item \imp{开放性问题}——不限于考查学生对知识的知晓情况，也可以考查学生应用知识解决问题的能力。
\end{itemize}
\end{keybox}

\subsection{RE-AIM评价框架}

\begin{keybox}
\textbf{RE-AIM框架：}于1999年由美国学者Russell E. Glasgow提出，目前已广泛地应用于公共卫生和项目评价的各个领域。由与健康行为干预有关的5个维度构成：

\begin{enumerate}[leftmargin=*]
    \item \imp{可及性（Reach）}——健康教育干预措施的人群覆盖情况，包括绝对数量、比例和代表性。在学校健康教育中，通常指学校健康教育课程的学生\imp{参与率}（实际参与人数/应参与人数）。参与率越高，可及性越高，实施效果越好。

    \item \imp{有效性（Effectiveness）}——健康教育干预措施对重要结局的影响，包括潜在的负面影响、生活质量和经济结局。在学校健康教育中，通常指学生完成一定阶段学习后，相关的\imp{健康知识、态度和行为}的改变情况。

    \item \imp{采纳性（Adoption）}——健康教育干预措施的组织机构人员参与和环境支持情况。在学校健康教育中，包括实际课时总量、健康教育课程师生比、健康教育专职或兼职教师的数量与质量、健康教育教学设备等。采纳性指标需要保持在合理适当的范围内，并非越高越好，否则会造成健康教育资源的浪费。

    \item \imp{实施性（Implementation）}——健康教育干预主体遵照方案各个要素执行的程度，包括按预期方案执行的一致性、所需时间、所做的调整和实施的成本。在学校健康教育中，可通过\imp{开课率}、实际执行过程中参与课程的师生比、健康教育课程\imp{目标达成率}等方面评估。由于学校健康教育的实施具有一定的时间跨度，实施性的评价也应当是\imp{动态的、过程性的}。

    \item \imp{持续性（Maintenance）}——健康教育干预措施的长期影响，包括维持、中断或适应的原因。在学校健康教育中，可通过两个方面评估：(1)项目实施一段时间后的可及性和有效性；(2)学校近几年健康教育课程的采纳性和实施性。
\end{enumerate}

\textbf{RE-AIM框架的意义：}通过五个维度的综合评价，不仅可以了解项目"是否有效"，还能评估其"在多大范围内有效""在什么条件下有效""能否持续推广"，弥补了传统评价仅关注有效性的不足，是当前国际上学校健康促进项目评价的推荐框架。
\end{keybox}

\newpage
